% META: {"id": "1NSI-ADGK-CO-005", "chapitre": "1NSI-ALGO-DICHO-GLOUTON-KNN", "type_objet": "corrige", "exercice_ref": "1NSI-ADGK-EX-005", "capacites": ["P-ALGO-05"], "sources_inspiration": [], "status": "verified"}
% BEGIN-VERIFY
% def sac_a_dos_glouton(objets, capacite):
%     objets_tries = sorted(objets, key=lambda o: o[2] / o[1], reverse=True)
%     charge = 0
%     valeur_totale = 0
%     choisis = []
%     for nom, poids, valeur in objets_tries:
%         if charge + poids <= capacite:
%             choisis.append(nom)
%             charge += poids
%             valeur_totale += valeur
%     return choisis, valeur_totale
% objets = [("X", 6, 30), ("Y", 5, 20), ("Z", 5, 20)]
% choisis, valeur = sac_a_dos_glouton(objets, 10)
% assert choisis == ["X"]
% assert valeur == 30
% assert 20 + 20 > valeur
% END-VERIFY
\begin{corrige}{1NSI-ADGK-EX-005}
\textbf{1.} Rapports valeur/poids : \lstinline{X} : $30/6=5$ ; \lstinline{Y} : $20/5=4$ ;
\lstinline{Z} : $20/5=4$. L'algorithme examine donc \lstinline{X} en premier (rapport le
plus élevé), puis \lstinline{Y} et \lstinline{Z}.

\textbf{2.} L'algorithme choisit d'abord \lstinline{X} (poids $6\leqslant10$, charge $=6$).
Il ne peut ensuite prendre ni \lstinline{Y} ni \lstinline{Z} (poids $5$ chacun, charge
$6+5=11>10$). Résultat : \lstinline{["X"]}, valeur totale $30$.

\textbf{3.} En prenant \lstinline{Y} et \lstinline{Z} ensemble (poids $5+5=10\leqslant10$),
on obtient une valeur de $20+20=40>30$. L'algorithme glouton n'est donc \textbf{pas
optimal} dans ce cas : privilégier le meilleur rapport valeur/poids objet par objet ne
garantit pas la meilleure combinaison globale, exactement comme pour le rendu de monnaie
avec un système de pièces non standard.
\end{corrige}
